\section{A prioritized research program}

\subsection{First priority: make the K-theory bridge independently checkable}

The most concrete near-term goal is a standalone theorem with a standard
cohomological statement and an explicit finite-dimensional output. Its
input should be a nonzero ordinary strict Selmer group for the rational
negative Tate twist. Its output should name all admissible witness
types, local levels, coefficient primes, and the trace congruences
at the probes actually used by the calculation.

There are two coherent choices:

\begin{enumerate}
\item \textbf{Retain pointwise witnesses.} At each probe, include
every admissible branch in the determinant product. Use the same
strict comparison for that probe and take a gcd only after those
products are formed. This is the route closest to the original
numerical bound.
\item \textbf{Extract a common representation.} Prove the stronger
construction-level statement used in the conditional Pintar route,
then import the precise Bernoulli-unit branch exclusions. This
can reduce the Hecke spaces, but adds a genuine theorem to the
dependency chain.
\end{enumerate}

The exact screening result in Theorem~\ref{thm:screening} is ready
for the second route at degrees $124$ and $128$. It removes the
specific problem of finding a two-field unit-compatible choice
for a large coefficient prime. It does not identify which
Hecke eigenvalues survive, prove their target congruences, or
show that the necessary determinant is a unit.

A concrete completion plan for each of those degrees is:

\begin{enumerate}
\item Fix the original strict-Selmer theorem, its local hypotheses,
and the common versus pointwise quantifiers.
\item Decide whether to retain only odd-discriminant fields or
validate the Gaussian modification. Keep the source condition
explicit in the theorem statement.
\item Generate the required integral Hecke operators and verify
the normalization of geometric Frobenius. For the inert Eisenstein
probe the trace is $(T_2+3I)/4$, not $T_2$.
\item Produce exact determinant or modular-inverse certificates
covering every remaining coefficient prime. Give B\'ezout
identities or factorizations with primality certificates where
large factors must be identified.
\item Combine the finite certificates with the standard small-prime
and $2$-primary comparisons, then check the complete dependency
closure of the resulting vanishing theorem.
\end{enumerate}

The arithmetic sizes now have a measured scale: the first-field
negative ranks are $550$ and $605$, and the second-field cost
depends strongly on the chosen route. Extending the Bernoulli table
farther is cheap, but it does not resolve the expensive Hecke or
cohomological parts. The next meaningful milestone is a certified
new determinant together with its exact interpretation.

\subsection{Second priority: the nonabelian cubic in mean}

The conductor-$23$ class-character target has several advantages.
Its analytic continuation, functional equation, and factorization
are already available. There is one genuinely new degree-two
factor. Character orthogonality gives an exact positive
principal-ideal selector. The Mellin residues have the correct
sign without a zero-avoidance family.

The missing arithmetic theorem can now be stated in several
forms. The positive-part square bound
\[
 \int_X^{2X}(\mathcal R_H(x)_+)^2\,dx
       \ll_\varepsilon X^\varepsilon
\]
is particularly compatible with a proof that already uses a
second moment for spanning. This compatibility is a research
suggestion, not an implication: the source's second moment is
over forms and translates, whereas this target is over the
prime-sum scale $x$. A bridge between those averaging variables
would have to be proved.

The rank-sensitive formula provides a separate concrete target:
control the mean of
\[
 U+\frac{L+2I}{\alpha},
\]
after the prime statistic and amplification have been normalized.
Average control of $L,I,\alpha$ separately is insufficient.
Successful control of this weighted loss at square-root
amplification would supply the positive-part criterion.

Alternatively, a construction that yields a signed moving-box
estimate directly can exploit Theorem~\ref{thm:moving-window}.
If the raw determinant inequality has a varying coefficient in
front of $\mathcal R_H$, one cannot discard that coefficient
merely because it is bounded: the remainder changes sign.
Either normalize before averaging or prove a zero-free
Mellin multiplier for the exact weight.

\subsection{Other prominent conjectures}

\begin{longtable}{@{}P{0.25\linewidth}P{0.43\linewidth}P{0.25\linewidth}@{}}
\toprule
Target & Required additional content & Assessment\\
\midrule
\endhead
Hecke RH over a quadratic field &
An ideal-class-sensitive arithmetic construction with the correct
coefficient identity and averaged or pointwise loss bounds. &
Specific first target isolated; no new zero bound proved.\\[5pt]
General Dedekind/automorphic RH &
New base fields, local types, integrality and spanning estimates;
where needed, analytic continuation and noncancellation. &
Plausible broad program, substantially beyond replacing names.\\[5pt]
Vandiver in all degrees &
Control of the appropriate even cyclotomic components uniformly
in the coefficient prime and weight. &
The source's odd-cocycle bound does not provide this.\\[5pt]
Artin's primitive-root density &
Zero estimates for the relevant Kummer number-field families,
with the needed uniformity as the extension varies. &
Dirichlet GRH alone does not match the required family.\\[5pt]
Zero simplicity, pair correlation, or linear independence &
Information about multiplicity and relative zero ordinates,
not only exclusion away from the critical line. &
No such information follows from the identified comparison.\\[5pt]
Twin primes, Goldbach, Elliott--Halberstam &
Correlation or average-over-moduli estimates exceeding
single-function zero location. &
No complete transfer supplied by this analysis.\\[5pt]
BSD &
A rank and leading-coefficient comparison for the relevant
elliptic-curve objects and their ordinary Selmer groups. &
Existence of favorable auxiliary twists is insufficient.\\
\bottomrule
\end{longtable}

These are statements about what the inspected method currently supplies.
They are not impossibility claims about future proofs. Even a successful
fixed-field nonabelian instance would not automatically give the
discriminant-uniform estimates needed in a varying-field application.

\subsection{What has actually been resolved in this investigation}

The completed mathematical results are the two averaged analytic criteria,
the mean criterion for combined rank and denominator losses, the elementary length-certificate
refinement with its exact hypothesis, and the finite Bernoulli and
representation-dimension theorems. The small Hecke operators have been
reconstructed independently from their published neighbor formula.

The investigation has not supplied the new arithmetic estimate needed
for the conductor-$23$ Hecke factor, completed the strict-Selmer
formalization, or extended the claimed K-group vanishing range.
Those boundaries are part of the result: they specify exactly what
would turn a method transfer into a solved open problem.
